% Draft text for the Overleaf report: Experimental Design and Evaluation.
% Source of truth for definitions: docs/evaluation.md. Update both together.

\section{Experimental Design}

\paragraph{Environment and episodes.}
Each episode is a seeded Minesweeper board rendered by our sandbox. The agent receives
only the window screenshot and issues normalized left or right clicks, which the sandbox
converts to pixels. The sandbox reads the true game state only to end the episode and to
grade it; this state is never shown to the model. An episode ends when the game is won or
lost, when the model emits \texttt{finish} or \texttt{fail}, on an unparseable or disallowed
action, or when the task's step or time limit is reached.

\paragraph{Harness and reproducibility.}
An experiment is one configuration (model, prompt, memory settings) run over every
benchmark task for $N$ trials. In trial $k$ the board seed is the task seed plus $k$, so
trials cover different boards while every configuration sees the same boards. Model
sampling is seeded per trial. For every experiment the harness records a manifest with the
full configuration, the git commit, the hash of every task file, and the platform. It also
saves the complete step log (raw model output, parsed action, execution result) and the
exact screenshot given to the model at each step. We report only experiments run from a
clean commit.

\paragraph{Baselines.}
We bracket the agent's performance with three reference policies:
(i) \emph{random click}, uniform over the window;
(ii) \emph{random cell}, which clicks a random unopened cell found from screenshot pixels;
and (iii) a \emph{perfect-perception solver}, which reads the visible board (revealed
numbers and flags, never the mine locations), applies single-cell deduction, and guesses
when stuck. The random policies are a floor for clicking without reasoning; the solver is
a ceiling for simple reasoning with perfect perception, and also confirms that each task
can be solved within its step limit. Our main baseline is the pretrained VLM with the same
prompt and no memory, compared against the full Vembody agent.

\paragraph{Ablations.}
Starting from the full agent, we disable one component of Budgeted Verified Transition
Memory at a time: transition verification, retained region crops, textual state memory,
and budgeted crop selection. All other settings, the tasks, and the seeds stay the same.

\section{Evaluation}

\paragraph{Metrics.}
The primary metric is \emph{success rate}, the fraction of episodes in which every safe
cell is revealed without revealing a mine. We expect an unmodified VLM to rarely win, so we
also report \emph{progress}: the fraction of safe cells revealed when the episode ends.
Efficiency is measured by steps per episode and wall-clock time. Action validity is the
\emph{invalid-action rate}, the fraction of model responses that fail parsing or
validation. We also count \emph{ineffective actions}, executed clicks that did not change
the board, and separately \emph{off-grid clicks}, which isolate grounding errors. A
\emph{setback} is an invalid or ineffective step. It counts as \emph{recovered} when the
next step changes the board, and \emph{recovery rate} is recoveries divided by setbacks.
Finally, we report median and 95th-percentile model latency per step and, where the model
runtime exposes them, prompt tokens and images per step as a measure of context usage.

\paragraph{Statistics.}
Success rates are reported with 95\% Wilson score intervals, which remain valid for small
samples and rates near 0 or 1. Mean progress is reported with 95\% percentile-bootstrap
intervals (2{,}000 resamples). Action-level rates are pooled over all steps in a
configuration. Comparisons between configurations use the same tasks and seeds, so we
also report per-task results to show where differences arise.

\paragraph{Reference results.}
On the beginner board ($9{\times}9$, 10 mines, 20 seeds, 100-step limit) the random-cell
policy wins 0\% of episodes (95\% CI 0--16\%) with mean progress 49\%, and the
perfect-perception solver wins 75\% (53--89\%) with mean progress 85\%. Random progress is
substantial because the first click is always safe and often opens a large region, so we
compare agent progress against this floor rather than against zero.
